\begingroup
\linespread{1}\selectfont
\renewcommand{\arraystretch}{0.95}
\centering
\tikzset{
  >=Latex,
  frame/.style={draw=black, rounded corners=3pt, line width=0.65pt, fill=white},
  header/.style={draw=black, rounded corners=2pt, line width=0.65pt, fill=black!8,
    align=center, font=\thesisfigurefont\bfseries, minimum height=7mm},
  data/.style={draw=black, rounded corners=2pt, line width=0.55pt, fill=black!3,
    align=center, font=\thesisfigurefont, minimum height=7mm, inner xsep=4pt},
  process/.style={draw=black, rounded corners=2pt, line width=0.55pt, fill=white,
    align=center, font=\thesisfigurefont, minimum height=7mm, inner xsep=4pt},
  note/.style={align=center, font=\thesisfigurefont, text=black},
  flow/.style={-{Latex[length=1.7mm,width=1.1mm]}, draw=black, line width=0.55pt},
  exchange/.style={<->, draw=black, line width=0.55pt}
}

\resizebox{\thesisfigurewidth}{!}{%
\begin{tikzpicture}[x=1.65cm,y=0.74cm]

\begin{scope}[shift={(0,16.35)}]
  \draw[frame] (0,0) rectangle (8.35,4.75);
  \node[header, minimum width=7.75cm] at (4.175,4.28)
    {参照：常规集中式离线监督};
  \node[data, minimum width=2.05cm] (s-train) at (1.45,2.92)
    {预先汇集的\\训练数据};
  \node[process, minimum width=1.70cm] (s-opt) at (4.05,2.92)
    {一次性优化};
  \node[process, minimum width=1.65cm] (s-model) at (6.75,2.92)
    {固定模型};
  \node[data, minimum width=2.05cm] (s-test) at (1.45,1.13)
    {独立测试样本\\相对稳定条件};
  \node[process, minimum width=1.70cm] (s-pred) at (6.75,1.13)
    {预测与评价};
  \draw[flow] (s-train) -- (s-opt);
  \draw[flow] (s-opt) -- (s-model);
  \draw[flow] (s-test) -- (4.05,1.13) -- (s-pred);
  \draw[flow] (s-model) -- (s-pred);
\end{scope}

\begin{scope}[shift={(0,10.90)}]
  \draw[frame] (0,0) rectangle (8.35,4.75);
  \node[header, minimum width=7.75cm] at (4.175,4.28)
    {元学习：任务组织与快速适应};
  \node[data, minimum width=1.95cm] (m-source) at (1.30,2.93)
    {多个源任务\\$\mathcal{T}_1,\ldots,\mathcal{T}_K$};
  \node[process, minimum width=1.65cm] (m-meta) at (4.05,2.93)
    {元训练};
  \node[process, minimum width=1.80cm] (m-init) at (6.80,2.93)
    {可适应初始化};
  \node[data, minimum width=1.95cm] (m-support) at (1.30,1.13)
    {新任务少量数据\\支持集};
  \node[process, minimum width=1.65cm] (m-adapt) at (4.05,1.13)
    {任务内适应};
  \node[process, minimum width=1.80cm] (m-task) at (6.80,1.13)
    {新任务模型};
  \draw[flow] (m-source) -- (m-meta);
  \draw[flow] (m-meta) -- (m-init);
  \draw[flow] (m-support) -- (m-adapt);
  \draw[flow] (m-init.south) -- ++(0,-.4) -| (m-adapt.north);
  \draw[flow] (m-adapt) -- (m-task);
\end{scope}

\begin{scope}[shift={(0,5.45)}]
  \draw[frame] (0,0) rectangle (8.35,4.75);
  \node[header, minimum width=7.75cm] at (4.175,4.28)
    {持续学习：时间到达与历史访问};
  \node[data, minimum width=1.15cm] (c-t1) at (0.95,2.85) {$\mathcal{T}_1$};
  \node[process, minimum width=1.15cm] (c-m1) at (2.45,2.85) {$\theta_1$};
  \node[data, minimum width=1.15cm] (c-t2) at (3.95,2.85) {$\mathcal{T}_2$};
  \node[process, minimum width=1.15cm] (c-m2) at (5.45,2.85) {$\theta_2$};
  \node[process, minimum width=1.15cm] (c-mt) at (7.60,2.85) {$\theta_t$};
  \draw[flow] (c-t1.east) -- (c-m1.west);
  \draw[flow] (c-m1.east) -- (c-t2.west);
  \draw[flow] (c-t2.east) -- (c-m2.west);
  \draw[flow] (c-m2.east) -- (6.20,2.85);
  \node[note, inner xsep=0pt] at (6.48,2.85) {$\cdots$};
  % 箭头在省略号之后、theta_t 方框之前结束，避免箭头侵入方框。
  \draw[flow] (6.72,2.85) -- ([xshift=-0.8mm]c-mt.west);
  \node[data, minimum width=2.55cm] (c-history) at (2.05,1.35)
    {可访问的历史信息\\完整／有限／无回放};
  \node[process, minimum width=2.55cm] (c-eval) at (6.20,1.35)
    {持续评价\\旧任务、当前任务与未来任务};
  \draw[flow] (c-history) -- (4.15,1.35) -- (c-eval);
  \draw[flow] (c-mt.south) -- ++(0,-.35) -| (c-eval.north);

\end{scope}

\begin{scope}[shift={(0,0)}]
  \draw[frame] (0,0) rectangle (8.35,4.75);
  \node[header, minimum width=7.75cm] at (4.175,4.28)
    {联邦学习：数据位置与协同优化};
  \node[data, minimum width=3.7cm] (f-c1) at (1.55,3.10)
    {客户端 1：本地数据};
  \node[data, minimum width=3.7cm] (f-c2) at (1.55,2.03)
    {客户端 2：本地数据};
  \node[note] at (1.55,1.38) {$\vdots$};
  \node[data, minimum width=3.7cm] (f-ck) at (1.55,0.70)
    {客户端 $K$：本地数据};
  \node[process, minimum width=2.05cm, minimum height=2.2cm] (f-server) at (4.65,2.03)
    {服务器\\聚合与下发};
  \node[process, minimum width=1.95cm] (f-global) at (7.10,2.03)
    {共享模型};
  % 各客户端连接独立端口，三条双向线等长且水平平行。
  \draw[exchange] (f-c1.east) -- (f-server.west |- f-c1.east);
  \draw[exchange] (f-c2.east) -- (f-server.west);
  \draw[exchange] (f-ck.east) -- (f-server.west |- f-ck.east);
  \draw[exchange] (f-server) -- (f-global);

\end{scope}


\end{tikzpicture}%
}
\endgroup
